
\documentclass[11pt,a4paper]{article}
\usepackage[a4paper,margin=2.05cm]{geometry}
\usepackage[spanish,es-noquoting]{babel}
\usepackage[T1]{fontenc}
\usepackage[utf8]{inputenc}
\usepackage{lmodern,microtype}
\usepackage{amsmath,amssymb,amsthm,mathtools}
\usepackage{booktabs,longtable,array,enumitem}
\usepackage{xcolor,hyperref}
\pagecolor{white}\color{black}
\setlength{\parindent}{0pt}
\setlength{\parskip}{0.65em}
\title{\bfseries N69 --- Ley de transición o axioma cilíndrico\\
\large Calendario mecánico, no-go afín y protocolo coinductivo de frontera}
\author{HMT / auditoría técnica}
\date{}
\begin{document}
\maketitle

\begin{abstract}
N69 ataca el blanco posterior a N68: convertir el protocolo de frontera en una ley de transición autónoma. Se prueba que el calendario abrir/podar es una ley mecánica exacta dada por \(K(t)=\lfloor6(t+1)\log_{1000}3\rfloor\). Se prueban candidatos de transición afín para bloques y firmas, con y sin fase \(12/22\), y no se obtiene una ley robusta. El resultado disciplinado es: seguir intentando derivar \(\mathcal F_{\rm HMT}\), pero si no se logra debe formularse como axioma coinductivo de verificación, no como teorema de generación infinita.
\end{abstract}

\section{Calendario mecánico}

La cantidad de triadas decimales visibles tras \(t\) bloques es:
\[
K(t)=\left\lfloor6(t+1)\log_{1000}3\right\rfloor.
\]
Por tanto:
\[
\Delta K_t=K(t+1)-K(t)\in\{0,1\}.
\]
Los stutters son los \(t\) con:
\[
\Delta K_t=0.
\]
Esto es una palabra mecánica/Beatty. No es una excepción: es parte infinita del protocolo.

\section{Tests de transición}

Se probaron mapas afines sobre:
\[
B_t,
\qquad
(q,a,c,\operatorname{colw},\rho_W,r)_t,
\qquad
(q,a)_t.
\]
También se probaron versiones con fase \(12\) y \(22\). El resumen:
\[
\begin{array}{c|c|c|c|c|c}
\text{modo}&\text{fase}&\dim_{in}&\dim_{out}&N&\text{consistente}\\
\hline
B & 0 & 18 & 18 & 45 & False \\
B & 12 & 30 & 18 & 45 & False \\
B & 22 & 40 & 18 & 45 & False \\
signature & 0 & 30 & 30 & 45 & False \\
signature & 12 & 42 & 30 & 45 & False \\
signature & 22 & 52 & 30 & 45 & True \\
qa & 0 & 12 & 12 & 45 & False \\
qa & 12 & 24 & 12 & 45 & False \\
qa & 22 & 34 & 12 & 45 & False \\
\end{array}
\]
La conclusión es:
\[
\boxed{
\text{no aparece una ley afín autónoma robusta de transición.}
}
\]

\section{Protocolo coinductivo}

La forma correcta del protocolo es:
\[
\boxed{
\text{Cyl}+\text{Beatty}+\text{Frontier}+\text{Survival}.
}
\]
Donde:
\[
\text{Cyl}=\text{compatibilidad por intersección de cilindros},
\]
\[
\text{Beatty}=K(t)=\lfloor6(t+1)\log_{1000}3\rfloor,
\]
\[
\text{Frontier}=B_t\mapsto(q,a,c,\operatorname{colw},\rho_W,r),
\]
\[
\text{Survival}=\text{rama admisible si sobrevive como cilindro inverso}.
\]

\section{Blanco exacto}

El teorema que falta, si queremos generación infinita completa, es:
\[
\boxed{
\mathcal F_{\rm HMT}(B_t,C_t)
=
(q,a,c,\operatorname{colw},\rho_W,r)_{t+1}.
}
\]
Si se deriva, la generación infinita queda en forma fuerte. Si no se deriva, la formulación honesta es:
\[
\boxed{
\text{axioma coinductivo de frontera y supervivencia.}
}
\]

\section{Dictamen}

\[
\boxed{
\text{verde: calendario mecánico y cilindro intervalar.}
}
\]
\[
\boxed{
\text{verde: protocolo de frontera + supervivencia como verificador.}
}
\]
\[
\boxed{
\text{rojo: transición afín simple.}
}
\]
\[
\boxed{
\text{amarillo: transición autónoma }\mathcal F_{\rm HMT}.
}
\]
\end{document}
